% Auto-generated by code/schema_mutation/phase10_analyze_nonobviousness.py
\begin{table}[t]
\centering
\caption{Non-obviousness control: extra reasoning without visibility does not provide the recovery channel, while making the evolved rule visible gives an upper-bound recovery signal.}
\label{tab:nonobviousness-control}
\small
\begin{tabular}{p{0.28\linewidth}rccp{0.20\linewidth}}
\toprule
Condition & N & Success & Hidden violation & Interpretation \\
\midrule
O0 + more reasoning & 96 & 3/96 (3.1%) & 92/96 (95.8%) & no visible rule \\
O0 + reflection & 95 & 6/95 (6.3%) & 88/95 (92.6%) & no visible rule \\
Rule-visible upper bound & 95 & 71/95 (74.7%) & 0/95 (0.0%) & rule visible \\
O0 standard reference & 96 & 0/96 (0.0%) & 96/96 (100.0%) & matched prior \\
\bottomrule
\end{tabular}
\vspace{0.3em}
\footnotesize{N counts ok rows only; one reflection cell failed and one rule-visible cell timed out, and these infrastructure rows are not counted as agent failures. The O0 standard reference row, when shown, is a matched prior frozen-result reference rather than a new Phase 10C call.}
\end{table}
